\documentclass[11pt,letterpaper]{article}
\usepackage[margin=1in]{geometry}
\pdfminorversion=6
\pdfobjcompresslevel=0
\usepackage{graphicx}
\usepackage{booktabs}
\usepackage{amsmath}
\usepackage{hyperref}
\usepackage{xcolor}
\usepackage{float}
\usepackage{enumitem}
\usepackage{multirow}

\title{Sunset Temperature Prediction at Rubin Observatory\\
\large Machine Learning Approaches for 3-Hour Forecasting}
\author{Analysis Report}
\date{January 2026}

\begin{document}

\maketitle

\begin{abstract}
We present a comprehensive analysis of algorithms for predicting exterior temperature at sunset at the Vera C.\ Rubin Observatory on Cerro Pach\'on, Chile. Using approximately two years of 15-minute interval temperature data (October 2023--December 2025), we developed and evaluated multiple prediction approaches including traditional methods and machine learning models. Models were \textbf{trained on odd calendar days} and \textbf{tested on even calendar days} to ensure rigorous out-of-sample evaluation. Our best model (Random Forest with 3-day lookback) achieves an \textbf{RMS error of 1.01$^\circ$C} at 3 hours before sunset with \textbf{bias of only $-$0.04$^\circ$C}, representing a 61\% improvement over the persistence baseline.
\end{abstract}

\section{Introduction}

Accurate prediction of sunset temperature is critical for telescope operations, particularly for thermal equilibration of optical components. This report documents our development of a prediction system with the following objectives:
\begin{itemize}[noitemsep]
    \item \textbf{Primary metric}: Minimize RMS error at 3 hours before sunset
    \item \textbf{Secondary goal}: Minimize systematic bias
    \item \textbf{Validation strategy}: Train on odd calendar days, test on even calendar days
\end{itemize}

\subsection{Data Overview}
\begin{itemize}[noitemsep]
    \item \textbf{Location}: Cerro Pach\'on, Chile (30$^\circ$14'40'' S, 70$^\circ$44'58'' W)
    \item \textbf{Time period}: October 2023 -- December 2025
    \item \textbf{Total valid days}: 771 days
    \item \textbf{Temporal resolution}: 15-minute intervals
    \item \textbf{Timestamps}: UTC (determined from diurnal cycle analysis)
\end{itemize}

\subsection{Train/Test Split Strategy}

To ensure rigorous out-of-sample evaluation, we split the data by calendar day number:
\begin{itemize}[noitemsep]
    \item \textbf{Training set (odd days)}: 394 days (days 1, 3, 5, 7, ... of each month)
    \item \textbf{Testing set (even days)}: 377 days (days 2, 4, 6, 8, ... of each month)
\end{itemize}

This interleaved split ensures that training and testing data span the same time periods and seasonal conditions, eliminating temporal bias while maintaining strict separation. Figure~\ref{fig:split} illustrates the data distribution.

\begin{figure}[H]
\centering
\includegraphics[width=\textwidth]{fig1_data_split.png}
\caption{\textbf{Left}: Visualization of odd/even day split across the dataset. \textbf{Right}: Sunset temperature distributions for training (odd days, blue) and testing (even days, orange) sets, showing similar statistical properties.}
\label{fig:split}
\end{figure}

\section{Development History}

\subsection{Phase 1: Baseline Models}
Initial implementation included:
\begin{enumerate}[noitemsep]
    \item \textbf{Persistence}: Current temperature as prediction (RMS = 2.57$^\circ$C at 3h)
    \item \textbf{Historical Mean}: Seasonal average ($\pm$15 day window) (RMS = 3.60$^\circ$C)
    \item \textbf{Linear Extrapolation}: 2-3 hour trend projection (RMS = 2.74--3.01$^\circ$C)
    \item \textbf{Fourier Decomposition}: 24h, 12h, 6h, 3h harmonics with polynomial baseline
\end{enumerate}

\textit{Finding}: Fourier methods underperformed persistence due to non-sinusoidal temperature evolution near sunset.

\subsection{Phase 2: Machine Learning Models}
Implemented ensemble and neural network approaches:
\begin{itemize}[noitemsep]
    \item Random Forest, Gradient Boosting, XGBoost, HistGradientBoosting
    \item Support Vector Regression (RBF kernel)
    \item Multi-layer Perceptron (MLP)
    \item Ridge regression
\end{itemize}

\textbf{Features}: Current and lookback temperatures (0.5--4h), temperature trends, daily statistics (max, min, range), and seasonal encoding (sin/cos of day-of-year).

\textbf{Result}: Random Forest achieved \textbf{1.08$^\circ$C RMS}---a 58\% improvement over persistence.

\subsection{Phase 3: Delta Prediction and Bias Correction}
Key improvements:
\begin{enumerate}[noitemsep]
    \item \textbf{Delta prediction}: Predict temperature \textit{change} from current to sunset, rather than absolute temperature
    \item \textbf{Calibration-based bias correction}: Use held-out calibration set (15\% of training data) to measure and subtract systematic bias
\end{enumerate}

\subsection{Phase 4: Lookback Optimization}
Investigated optimal historical data window. \textbf{Finding}: 3 days (72 hours) of lookback is optimal, as integral day boundaries correctly capture diurnal temperature cycles.

Final feature set includes: previous days' mean temperatures, temperature ranges, and sunset temperatures.

\section{Results}

\subsection{Performance Overview}

Figure~\ref{fig:rms_heatmap} shows RMS prediction error across all models and lead times.

\begin{figure}[H]
\centering
\includegraphics[width=0.9\textwidth]{fig2_rms_heatmap.png}
\caption{RMS prediction error (°C) by model and lead time, including both numerical and text-based approaches. Green indicates lower error. Random Forest achieves the best performance (1.06°C at 3h), outperforming text-based LSTM (1.31°C) and Transformer (1.43°C). White horizontal line separates numerical (above) from text-based (below) methods.}
\label{fig:rms_heatmap}
\end{figure}

\subsection{Multiple Error Metrics}

Figure~\ref{fig:median_mae} presents median and mean absolute error heatmaps.

\begin{figure}[H]
\centering
\includegraphics[width=\textwidth]{fig3_median_mae_heatmaps.png}
\caption{\textbf{Left}: Median absolute error. \textbf{Right}: Mean absolute error. The median is consistently lower than the mean, indicating a right-skewed error distribution with occasional large outliers.}
\label{fig:median_mae}
\end{figure}

\subsection{Error Distributions}

Figure~\ref{fig:histograms} shows prediction error histograms at 3 hours lead time.

\begin{figure}[H]
\centering
\includegraphics[width=\textwidth]{fig4_error_histograms.png}
\caption{Error distribution histograms at 3-hour lead time for all models tested on even days. Red vertical line indicates zero error; green dashed line shows mean bias. Random Forest shows the tightest distribution with minimal bias ($-$0.04°C).}
\label{fig:histograms}
\end{figure}

\subsection{Cumulative Distribution Functions}

Figure~\ref{fig:cdfs} presents cumulative probability distributions of absolute errors.

\begin{figure}[H]
\centering
\includegraphics[width=\textwidth]{fig5_cdfs.png}
\caption{\textbf{Left}: CDF for Random Forest at different lead times, showing 70\% of 3-hour predictions have error $<$1°C. \textbf{Right}: CDF comparison across all models at 3 hours.}
\label{fig:cdfs}
\end{figure}

\subsection{RMS vs Lead Time}

Figure~\ref{fig:rms_leadtime} shows how prediction error grows with lead time.

\begin{figure}[H]
\centering
\includegraphics[width=0.85\textwidth]{fig6_rms_vs_leadtime.png}
\caption{RMS error vs lead time for all models. Error grows approximately linearly at $\sim$0.15°C per hour. At 3 hours, best models achieve RMS $\approx$ 1.0°C.}
\label{fig:rms_leadtime}
\end{figure}

\subsection{Summary Table}

Table~\ref{tab:results} presents final performance metrics at 3 hours before sunset for all methods, including both numerical feature-based approaches and text-based embedding approaches.

\begin{table}[H]
\centering
\caption{Comprehensive Model Performance at 3-Hour Lead Time (Tested on Even Days)}
\label{tab:results}
\begin{tabular}{llccccc}
\toprule
\textbf{Type} & \textbf{Model} & \textbf{RMS (°C)} & \textbf{MAE (°C)} & \textbf{Bias (°C)} & \textbf{\% $<$ 1°C} \\
\midrule
\multirow{7}{*}{\rotatebox{90}{Numerical}}
& \textbf{RF + Sample Weight} & \textbf{1.047} & 0.812 & $+$0.113 & 68.7 \\
& Random Forest & 1.057 & 0.822 & $+$0.120 & 69.0 \\
& Ridge & 1.115 & 0.885 & $+$0.074 & 64.2 \\
& MLP & 1.119 & 0.884 & $+$0.152 & 62.6 \\
& GradientBoosting & 1.125 & 0.867 & $+$0.097 & 65.5 \\
& HistGradientBoosting & 1.127 & 0.881 & $+$0.128 & 65.8 \\
& XGBoost & 1.152 & 0.879 & $+$0.142 & 66.3 \\
\midrule
\multirow{2}{*}{\rotatebox{90}{Text}}
& LSTM+Attention & 1.31 & 1.04 & $+$0.41 & 54.9 \\
& Transformer & 1.43 & 1.15 & $+$0.25 & 52.3 \\
\midrule
& Persistence (baseline) & 2.53 & 2.27 & $+$2.21 & 14.3 \\
\bottomrule
\end{tabular}
\end{table}

\textbf{Key observations}:
\begin{itemize}[noitemsep]
    \item Random Forest with sample weighting achieves the best overall RMS (1.047°C)
    \item All numerical approaches outperform text-based approaches at 3 hours
    \item Text-based LSTM (1.31°C) is 25\% worse than best numerical; Transformer (1.43°C) is 37\% worse
    \item All ML methods substantially outperform persistence baseline (2.53°C)
\end{itemize}

\section{Error Analysis}

\subsection{Temporal Evolution}

Figure~\ref{fig:time_evolution} shows how prediction errors evolve over the two-year dataset.

\begin{figure}[H]
\centering
\includegraphics[width=\textwidth]{fig7_error_time_evolution.png}
\caption{Temporal evolution of prediction errors. \textbf{Top left}: Individual errors over time with 30-day rolling mean. \textbf{Top right}: 30-day rolling RMS showing stable performance. \textbf{Bottom left}: Absolute errors over time. \textbf{Bottom right}: Monthly RMS (red=summer, green=winter, blue=transition seasons).}
\label{fig:time_evolution}
\end{figure}

\subsection{Seasonal Effects}

Figure~\ref{fig:seasonal} analyzes performance differences across seasons.

\begin{figure}[H]
\centering
\includegraphics[width=\textwidth]{fig8_seasonal_effects.png}
\caption{Seasonal analysis. \textbf{Top left}: Error distribution by season. \textbf{Top right}: RMS and MAE by season. \textbf{Bottom left}: Bias by season---Fall shows largest negative bias ($-$0.29°C), Spring shows positive bias (+0.16°C). \textbf{Bottom right}: Error vs day of year with 15-day binned RMS.}
\label{fig:seasonal}
\end{figure}

Table~\ref{tab:seasonal} summarizes seasonal performance metrics.

\begin{table}[H]
\centering
\caption{Seasonal Performance Summary}
\label{tab:seasonal}
\begin{tabular}{lcccc}
\toprule
\textbf{Season} & \textbf{RMS (°C)} & \textbf{MAE (°C)} & \textbf{Bias (°C)} & \textbf{N days} \\
\midrule
Summer (Dec--Feb) & 1.024 & 0.788 & $-$0.092 & 88 \\
Fall (Mar--May) & 1.026 & 0.810 & $-$0.293 & 90 \\
Winter (Jun--Aug) & 1.000 & 0.782 & $+$0.024 & 90 \\
Spring (Sep--Nov) & 0.988 & 0.786 & $+$0.162 & 109 \\
\bottomrule
\end{tabular}
\end{table}

\textbf{Key observation}: Seasonal bias varies from $-$0.29°C (Fall) to $+$0.16°C (Spring). This suggests that season-specific bias correction could further improve performance.

\subsection{Dependence on Temperature Trends}

Figure~\ref{fig:trends} examines whether errors depend on multi-day temperature trends.

\begin{figure}[H]
\centering
\includegraphics[width=\textwidth]{fig9_trend_dependence.png}
\caption{Error dependence on temperature trends. \textbf{Top left}: Error vs 3-day trend. \textbf{Top right}: Error vs 7-day trend. \textbf{Bottom left}: RMS by trend category---stable conditions ($\pm$0.1°C/day) show lowest RMS (0.93°C). \textbf{Bottom right}: Error vs daily temperature range.}
\label{fig:trends}
\end{figure}

\textbf{Key findings}:
\begin{itemize}[noitemsep]
    \item Correlations between temperature trends and errors are weak ($r < 0.1$)
    \item \textbf{Stable days} (trend $\pm$0.1°C/day) show lowest RMS: 0.93°C
    \item \textbf{Strong cooling} days ($<-$0.5°C/day) show highest RMS: 1.09°C
    \item Days with larger diurnal temperature range show slightly higher errors ($r = 0.10$)
\end{itemize}

\section{Comprehensive Comparison: Numerical vs.\ Text-Based Approaches}

This section provides a detailed comparison of two fundamentally different philosophies for predicting sunset temperature: \textbf{numerical/feature-based approaches} (like Random Forest) and \textbf{text-based embedding approaches} (inspired by Large Language Models). Understanding their differences illuminates why each succeeds and where each has limitations.

\subsection{Philosophical Overview}

\begin{figure}[H]
\centering
\includegraphics[width=\textwidth]{fig20_philosophical_comparison.png}
\caption{Philosophical comparison between numerical and text-based approaches. \textbf{Top}: The numerical approach treats temperatures as floating-point numbers and engineers explicit features (means, trends, ranges) that domain expertise suggests are predictive. \textbf{Bottom}: The text-based approach treats temperatures as discrete ``words'' (ASCII strings like ``+12.5''), learns vector embeddings from data, and uses sequence models to discover patterns.}
\label{fig:philosophy}
\end{figure}

Table~\ref{tab:philosophy} summarizes the key philosophical differences.

\begin{table}[H]
\centering
\caption{Philosophical Differences: Numerical vs.\ Text-Based Approaches}
\label{tab:philosophy}
\begin{tabular}{p{3.5cm}p{5.5cm}p{5.5cm}}
\toprule
\textbf{Aspect} & \textbf{Numerical (Random Forest)} & \textbf{Text-Based (LSTM/Transformer)} \\
\midrule
\textbf{Data representation} & Floating-point numbers & Discrete ASCII strings (``+12.5'') \\
\textbf{Feature engineering} & Explicit: means, trends, ranges & Implicit: learned from data \\
\textbf{Domain knowledge} & Encoded in feature design & Discovered by model \\
\textbf{Model learns} & Decision boundaries in feature space & Embedding vectors + sequence patterns \\
\textbf{Interpretability} & Feature importance rankings & Attention weights, embedding similarity \\
\textbf{Data efficiency} & High (works with $<$1000 samples) & Lower (benefits from more data) \\
\textbf{Flexibility} & Requires new features for new tasks & Can adapt to new patterns automatically \\
\bottomrule
\end{tabular}
\end{table}

\subsection{Training Data Cutoff: What the Models ``See''}

A critical methodological point: for a 3-hour-ahead prediction, the model cannot see any data between the prediction time ($T-3$h) and sunset. This ensures the models are evaluated on true out-of-sample prediction.

\begin{figure}[H]
\centering
\includegraphics[width=\textwidth]{fig19_training_data_cutoff.png}
\caption{Training data cutoff illustrated for 3-hour, 2-hour, and 1-hour lead times. \textbf{Green regions}: Data available at prediction time (used for training/inference). \textbf{Red regions}: Data not available at prediction time (excluded from training). For 3-hour predictions, we drop approximately 12 temperature readings between $T-3$h and sunset. The gold star indicates the target sunset temperature.}
\label{fig:cutoff}
\end{figure}

Both numerical and text-based approaches use the same data cutoff:
\begin{itemize}[noitemsep]
    \item \textbf{3-hour prediction}: Uses data from $T-75$h to $T-3$h (3 days of lookback)
    \item \textbf{2-hour prediction}: Uses data from $T-74$h to $T-2$h
    \item \textbf{1-hour prediction}: Uses data from $T-73$h to $T-1$h
\end{itemize}

\section{Experimental: LLM-Style Embedding Approach}

As an experimental investigation, we explored whether treating temperature data as ``language'' could provide predictive power. This section describes an approach inspired by Large Language Models (LLMs), where we represent temperatures as ASCII text strings and process them using sequence models with learned embeddings.

\subsection{Motivation}

Modern language models like GPT achieve remarkable performance by:
\begin{enumerate}[noitemsep]
    \item Tokenizing text into discrete vocabulary elements
    \item Learning dense vector embeddings for each token
    \item Processing sequences with attention-based architectures
    \item Capturing long-range dependencies in sequential data
\end{enumerate}

We asked: \textit{Can we apply similar techniques to temperature time series?} Instead of predicting ``the next word,'' we use the sequence representation to predict a continuous temperature value.

\subsection{Step 1: Tokenization---Temperatures as ASCII Words}

The first step converts continuous temperature values into discrete ``words'' using ASCII string representation.

\begin{figure}[H]
\centering
\includegraphics[width=\textwidth]{fig12_tokenization.png}
\caption{The tokenization process. \textbf{(A)} Raw temperature time series at 15-minute intervals. \textbf{(B)} Quantization to 0.5°C precision reduces vocabulary size while preserving information. \textbf{(C)} Conversion to ASCII strings with explicit sign (e.g., $-5.23$°C $\rightarrow$ ``-5.0'', $+12.78$°C $\rightarrow$ ``+13.0''). \textbf{(D)} Vocabulary distribution showing 68 unique temperature ``words'' in our dataset.}
\label{fig:tokenization}
\end{figure}

The tokenization procedure:
\begin{enumerate}[noitemsep]
    \item \textbf{Quantize}: Round temperature to nearest 0.5°C
    \item \textbf{Format}: Convert to signed ASCII string: \texttt{f"\{temp:+.1f\}"}
    \item \textbf{Result}: Vocabulary of 68 unique ``words'' from ``-9.0'' to ``+23.5''
\end{enumerate}

This is analogous to how text tokenizers convert words to discrete token IDs, except our ``words'' are temperature strings like ``+12.5'' or ``-3.0''.

\subsection{Step 2: Sentence Construction from 3-Day Windows}

Each prediction instance becomes a ``sentence'' composed of all temperature readings from the past 3 days plus the current day up to the prediction time (3 hours before sunset).

\begin{figure}[H]
\centering
\includegraphics[width=\textwidth]{fig13_sentence_construction.png}
\caption{Construction of temperature ``sentences.'' \textbf{Top}: Three days of historical temperature data plus today's data up to $T-3$h forms the input window. The gold star marks the target sunset temperature we want to predict. \textbf{Bottom}: The resulting ``sentence'' of approximately 371 temperature words that serves as input to the model.}
\label{fig:sentence}
\end{figure}

A typical sentence contains $\sim$371 words (approximately 93 readings per day $\times$ 4 days). For example:
\begin{center}
\texttt{[``+8.0'', ``+8.5'', ``+9.0'', ``+9.5'', ..., ``+11.5'', ``+12.0'']}
\end{center}

This is directly analogous to feeding a sentence of 371 words to a language model.

\subsection{Step 3: Model Architecture}

\textbf{Key Insight}: Unlike GPT which predicts the probability distribution over the next word, we use the sequence encoding to predict a \textit{continuous} temperature value. This is a regression task, not a classification task.

\begin{figure}[H]
\centering
\includegraphics[width=\textwidth]{fig14_model_architecture.png}
\caption{Model architecture for the ASCII word embedding approach. Temperature ``words'' are embedded into 64-dimensional vectors, processed by a bidirectional LSTM with attention, concatenated with the current temperature, and passed through fully-connected layers to predict the temperature change $\Delta T$. The final prediction is $T_\text{sunset} = T_\text{current} + \Delta T$.}
\label{fig:architecture}
\end{figure}

Figure~\ref{fig:architecture_detail} provides a more detailed view of the training pipeline.

\begin{figure}[H]
\centering
\includegraphics[width=\textwidth]{fig22_text_model_architecture.png}
\caption{Detailed training pipeline for text-based models. \textbf{Step 1}: Temperature readings are tokenized into ASCII strings. \textbf{Step 2}: Each token is mapped to a learned 64-dimensional embedding vector. \textbf{Step 3}: The sequence is processed by either LSTM (reads sequentially) or Transformer (reads all at once with attention). \textbf{Step 4}: The context vector summarizes the entire temperature history. \textbf{Step 5}: A regression head predicts $\Delta T$, which is added to the current temperature. The key insight: we predict a \textit{number}, not the ``next word.''}
\label{fig:architecture_detail}
\end{figure}

We implemented two architectures:

\textbf{1. Bidirectional LSTM with Attention (Long Short-Term Memory):}
\begin{itemize}[noitemsep]
    \item Word embedding layer: 68 words $\rightarrow$ 64-dimensional vectors
    \item 2-layer bidirectional LSTM (128 hidden units)
    \item LSTM reads the sequence step-by-step, maintaining a ``memory'' of previous temperatures
    \item Attention mechanism weights the importance of each time step
    \item Concatenate context vector with current temperature
    \item Fully-connected layers: 257 $\rightarrow$ 64 $\rightarrow$ 32 $\rightarrow$ 1
\end{itemize}

\textbf{2. Transformer Encoder:}
\begin{itemize}[noitemsep]
    \item Word + positional embeddings (64-dimensional)
    \item 3-layer Transformer encoder (4 attention heads)
    \item Transformer reads all temperatures simultaneously via self-attention
    \item Mean pooling over sequence to create context vector
    \item Same regression head as LSTM
\end{itemize}

\textbf{Note on terminology}: LSTM is \textit{not} a Large Language Model. ``LLM'' refers specifically to large Transformer-based models like GPT and Claude. LSTM is an older recurrent architecture from the 1990s. We tested both to compare different sequence modeling approaches.

\subsection{Step 4: The Prediction Mechanism}

Figure~\ref{fig:prediction_mechanism} clarifies how prediction works in our approach versus traditional language models.

\begin{figure}[H]
\centering
\includegraphics[width=\textwidth]{fig16_prediction_mechanism.png}
\caption{The prediction mechanism explained. \textbf{(A)} Comparison with traditional LLMs: we predict a continuous temperature change, not a probability over vocabulary. \textbf{(B)} The model sees 3 days of history up to the prediction time (blue shaded region) and must predict the sunset temperature (gold star). \textbf{(C)} Delta prediction: the model predicts the temperature \textit{change} from current to sunset, then adds it to the current temperature. \textbf{(D)} Rationale for delta prediction.}
\label{fig:prediction_mechanism}
\end{figure}

The model predicts $\Delta T = T_\text{sunset} - T_\text{current}$, then computes:
\begin{equation}
\hat{T}_\text{sunset} = T_\text{current} + \Delta T_\text{predicted}
\end{equation}

This delta prediction strategy:
\begin{itemize}[noitemsep]
    \item Reduces target range: $\Delta T \in [-8, +4]$°C vs $T \in [-9, +24]$°C
    \item Simplifies the learning task: ``how much will temperature change?''
    \item Incorporates the current temperature, our most predictive feature
    \item Matches the approach used by our successful Random Forest model
\end{itemize}

\subsection{Step 5: Learned Embeddings}

The model learns 64-dimensional embeddings for each temperature ``word.'' Figure~\ref{fig:embeddings} visualizes these learned representations.

\begin{figure}[H]
\centering
\includegraphics[width=\textwidth]{fig15_embedding_space.png}
\caption{Analysis of learned temperature word embeddings. \textbf{(A)} PCA projection of the 64-dimensional embeddings shows that similar temperatures cluster together---the model has learned that ``+10.0'' and ``+10.5'' should have similar representations. \textbf{(B)} Cosine similarity matrix confirms that nearby temperatures have higher similarity (red diagonal band) while distant temperatures have lower similarity.}
\label{fig:embeddings}
\end{figure}

The embedding analysis reveals:
\begin{itemize}[noitemsep]
    \item Adjacent temperature words ($\Delta = 0.5$°C) have positive similarity
    \item Distant temperature words ($\Delta = 5$°C) have near-zero or negative similarity
    \item The model correctly learns the ordinal structure of temperatures
\end{itemize}

This is analogous to how word embeddings in NLP learn that ``king'' and ``queen'' are more similar than ``king'' and ``banana.''

\subsection{Results: LLM Approach Performance}

We evaluated the LLM approach with two refinements: (1) increased precision from 0.5°C to 0.1°C (vocabulary of 323 words), and (2) multiple lead times (3h, 2h, 1h before sunset).

\begin{figure}[H]
\centering
\includegraphics[width=\textwidth]{fig18_llm_expanded_results.png}
\caption{Comparison of Random Forest vs text-based models across lead times. \textbf{(A)} RMS error vs lead time---all models improve at shorter lead times, but Random Forest (green) consistently outperforms LSTM (blue) and Transformer (coral). \textbf{(B)} Percentage worse than Random Forest at each lead time---text-based models are 20--35\% worse at every lead time. \textbf{(C)} Bias comparison---Random Forest maintains near-zero bias while text-based models show positive bias. \textbf{(D)} Percentage of predictions within 1°C---Random Forest achieves 85\% at 1h vs 76\% for Transformer.}
\label{fig:llm_expanded}
\end{figure}

Table~\ref{tab:full_comparison} (in Section~\ref{sec:complete_comparison}) provides a comprehensive comparison of all three approaches across lead times.

\textbf{Key findings from expanded analysis:}
\begin{itemize}[noitemsep]
    \item \textbf{Random Forest wins at every lead time}: RF outperforms text-based models by 18--42\% at all lead times
    \item At \textbf{1 hour}: RF achieves 0.74°C vs Transformer 0.87°C (+18\%) and LSTM 0.93°C (+26\%)
    \item At \textbf{2 hours}: RF achieves 0.95°C vs LSTM 1.14°C (+20\%) and Transformer 1.22°C (+28\%)
    \item At \textbf{3 hours}: RF achieves 1.01°C vs LSTM 1.31°C (+30\%) and Transformer 1.43°C (+42\%)
    \item Increasing precision from 0.5°C to 0.1°C (323 vs 68 vocabulary words) provides \textbf{no significant improvement}
\end{itemize}

\subsection{Tokenization Granularity: 0.5°C vs.\ 0.1°C}

An important question: \textit{does finer temperature resolution improve predictions?} We tested both 0.5°C and 0.1°C quantization to understand how tokenization precision affects model performance.

\begin{figure}[H]
\centering
\includegraphics[width=\textwidth]{fig23_granularity_comparison.png}
\caption{Effect of tokenization granularity on text-based model performance. \textbf{(A)} Vocabulary sizes differ dramatically: 68 words at 0.5°C vs.\ 323 words at 0.1°C. \textbf{(B)} RMS comparison at 3-hour lead time shows no meaningful difference. \textbf{(C)} Example tokenization: the same temperature sequence produces different token strings at each precision. \textbf{(D)} Implications: finer granularity does not help---the prediction bottleneck is not tokenization precision.}
\label{fig:granularity}
\end{figure}

\begin{table}[H]
\centering
\caption{Tokenization Granularity Comparison at 3-Hour Lead Time}
\label{tab:granularity}
\begin{tabular}{lcccc}
\toprule
\textbf{Precision} & \textbf{Vocabulary} & \textbf{LSTM RMS} & \textbf{Transformer RMS} & \textbf{Difference} \\
\midrule
0.5°C & 68 words & 1.31°C & 1.41°C & --- \\
0.1°C & 323 words & 1.31°C & 1.43°C & $<$0.02°C \\
\bottomrule
\end{tabular}
\end{table}

\textbf{Key insight}: Increasing precision from 0.5°C to 0.1°C (4.75$\times$ larger vocabulary) provides \textit{no improvement}. This reveals that:
\begin{itemize}[noitemsep]
    \item The model already captures relevant patterns at 0.5°C resolution
    \item The prediction bottleneck is inherent unpredictability, not input precision
    \item Larger vocabulary means sparser training signal per token, which may hurt generalization
\end{itemize}

\subsection{Complete Results Comparison}
\label{sec:complete_comparison}

Figure~\ref{fig:complete} provides a comprehensive comparison of all approaches.

\begin{figure}[H]
\centering
\includegraphics[width=\textwidth]{fig21_complete_comparison.png}
\caption{Complete comparison of all prediction methods. \textbf{(A)} Bar chart comparing RMS at 3-hour lead time: Random Forest (1.01°C) outperforms text-based LSTM (1.31°C) and Transformer (1.43°C). \textbf{(B)} Comparison vs Random Forest benchmark. \textbf{(C)} RMS at 3 hours---Random Forest is the clear winner. \textbf{(D)} Summary: Random Forest wins at all lead times (0.74°C at 1h, 0.95°C at 2h, 1.01°C at 3h).}
\label{fig:complete}
\end{figure}

Table~\ref{tab:full_comparison} provides a comprehensive comparison of Random Forest, LSTM, and Transformer performance across all lead times.

\begin{table}[H]
\centering
\caption{Complete Model Comparison Across All Lead Times}
\label{tab:full_comparison}
\begin{tabular}{llcccc}
\toprule
\textbf{Lead Time} & \textbf{Model} & \textbf{RMS (°C)} & \textbf{MAE (°C)} & \textbf{Bias (°C)} & \textbf{\% $<$ 1°C} \\
\midrule
\multirow{3}{*}{3 hours} & \textbf{Random Forest} & \textbf{1.01} & \textbf{0.79} & $-$0.04 & \textbf{69.8} \\
 & LSTM+Attention & 1.31 & 1.04 & $+$0.41 & 54.9 \\
 & Transformer & 1.43 & 1.15 & $+$0.25 & 52.3 \\
\midrule
\multirow{3}{*}{2 hours} & \textbf{Random Forest} & \textbf{0.95} & \textbf{0.75} & $+$0.04 & \textbf{73.2} \\
 & LSTM+Attention & 1.14 & 0.91 & $+$0.19 & 59.4 \\
 & Transformer & 1.22 & 0.99 & $+$0.19 & 58.4 \\
\midrule
\multirow{3}{*}{1 hour} & \textbf{Random Forest} & \textbf{0.74} & \textbf{0.59} & $+$0.09 & \textbf{81.4} \\
 & LSTM+Attention & 0.93 & 0.72 & $+$0.12 & 74.5 \\
 & Transformer & 0.87 & 0.69 & $+$0.05 & 75.9 \\
\bottomrule
\end{tabular}
\end{table}

\textbf{Key observations from Table~\ref{tab:full_comparison}:}
\begin{itemize}[noitemsep]
    \item \textbf{Random Forest wins at ALL lead times}, not just at 3 hours
    \item At 1 hour, Random Forest (0.74°C) significantly outperforms both Transformer (0.87°C) and LSTM (0.93°C)
    \item The gap between Random Forest and text-based models narrows at shorter lead times
    \item Text-based models have consistently positive bias; Random Forest has near-zero bias
\end{itemize}

\subsection{Discussion: Comparing Text-Based and Numerical Approaches}

At the primary 3-hour lead time, Random Forest (1.01°C) significantly outperforms the text-based approaches: LSTM (1.31°C, 30\% worse) and Transformer (1.43°C, 42\% worse). The text-based models do learn meaningful patterns---they substantially beat persistence---but they cannot match the performance of well-engineered numerical features.

\textbf{Why does Random Forest win at 3 hours?}
\begin{itemize}[noitemsep]
    \item Explicit feature engineering encodes domain knowledge: trends, ranges, daily patterns
    \item Tree-based models excel at capturing non-linear relationships with limited data
    \item The hand-crafted features directly represent what matters for prediction
    \item $\sim$400 training samples is sufficient for feature-based models but limiting for neural networks
\end{itemize}

\textbf{Why does Transformer beat LSTM at short lead times?}
\begin{itemize}[noitemsep]
    \item Attention mechanism can directly focus on recent temperatures
    \item At 1h, the most relevant information is in the last few readings
    \item LSTM's sequential processing may introduce unnecessary complexity
    \item At 1h, Transformer (0.87°C) beats LSTM (0.93°C) but still loses to Random Forest (0.74°C)
\end{itemize}

\textbf{When might text-based approaches excel?}
\begin{itemize}[noitemsep]
    \item With significantly more training data (years to decades)
    \item When incorporating heterogeneous data (text weather reports, satellite imagery)
    \item For multi-site prediction where transfer learning could help
    \item With pre-trained foundation models for time series
    \item When feature engineering expertise is unavailable
\end{itemize}

\textbf{The fundamental trade-off:}
\begin{itemize}[noitemsep]
    \item \textbf{Numerical/feature-based}: Requires domain expertise to design features, but highly data-efficient and interpretable
    \item \textbf{Text-based/embedding}: Learns representations automatically, but requires more data and is less interpretable
\end{itemize}

\subsection{Conclusion: LLM Experiment}

Treating temperatures as ASCII ``words'' and using language model architectures is a creative approach that \textit{does work}---achieving 0.87°C RMS at 1 hour before sunset, with 76\% of predictions within 1°C. The models correctly learn that similar temperatures should have similar embeddings. However, for the primary 3-hour prediction task, Random Forest (1.01°C) still outperforms the LLM approach (1.31°C). This experiment demonstrates both the flexibility of modern deep learning and the continued importance of matching model architecture to problem structure.

\section{Additional Methodology Exploration}

Building on the success of Random Forest, we explored several additional methodologies to potentially improve prediction accuracy or better understand when predictions are likely to be less reliable. This section summarizes six additional analyses.

\subsection{Error Condition Analysis}

We investigated what conditions correlate with larger prediction errors. Understanding when predictions are less reliable enables better uncertainty quantification.

\begin{figure}[H]
\centering
\includegraphics[width=\textwidth]{fig24_error_analysis.png}
\caption{Analysis of conditions correlating with prediction errors. \textbf{(A)} Error vs temperature standard deviation over 3 days. \textbf{(B)} Error vs 2-hour rate of change---the strongest predictor. \textbf{(C)} Error vs 24-hour temperature change. \textbf{(D)} RMS by rate-of-change category showing that slow-changing conditions yield RMS=0.88°C while fast-changing conditions yield RMS=1.14°C.}
\label{fig:error_analysis}
\end{figure}

\textbf{Key finding}: The 2-hour rate of temperature change is the best predictor of error magnitude:
\begin{itemize}[noitemsep]
    \item \textbf{Slow change} ($|$rate$|<$0.15°C/h): RMS = 0.88°C
    \item \textbf{Fast change} ($|$rate$|\geq$0.15°C/h): RMS = 1.14°C
    \item This 30\% difference enables uncertainty quantification
\end{itemize}

\subsection{Fine-Grained Lead Time Analysis}

We analyzed prediction performance at 15-minute intervals from 0.25 to 4 hours before sunset.

\begin{figure}[H]
\centering
\includegraphics[width=\textwidth]{fig25_fine_grained_leadtime.png}
\caption{Fine-grained lead time analysis at 15-minute intervals. \textbf{(A)} RMS error grows smoothly from 0.30°C at 15 minutes to 1.07°C at 3 hours, crossing 1.0°C at approximately 2.75 hours. \textbf{(B)} MAE follows similar pattern. \textbf{(C)} Percentage of predictions within 1°C decreases from 99.2\% at 15 min to 66.8\% at 3 hours. \textbf{(D)} Random Forest improvement over persistence.}
\label{fig:fine_grained}
\end{figure}

Table~\ref{tab:finegrained} summarizes performance at key lead times.

\begin{table}[H]
\centering
\caption{Performance at Fine-Grained Lead Times}
\label{tab:finegrained}
\begin{tabular}{lccc}
\toprule
\textbf{Lead Time} & \textbf{RMS (°C)} & \textbf{MAE (°C)} & \textbf{\% $<$ 1°C} \\
\midrule
15 min & 0.30 & 0.23 & 99.2 \\
30 min & 0.44 & 0.34 & 97.9 \\
1 hour & 0.72 & 0.56 & 86.7 \\
2 hours & 0.95 & 0.74 & 74.0 \\
3 hours & 1.07 & 0.84 & 66.8 \\
\bottomrule
\end{tabular}
\end{table}

\textbf{Key finding}: Error grows at approximately \textbf{0.177°C per hour}. The 1.0°C RMS threshold is crossed at approximately 2.75 hours before sunset.

\subsection{Exponentially Weighted Lookback}

Instead of a hard 3-day cutoff, we tested exponentially decaying weights: $w = e^{-\lambda \cdot \text{days\_ago}}$

\begin{figure}[H]
\centering
\includegraphics[width=\textwidth]{fig27_exponential_weighting.png}
\caption{Exponential weighting analysis. \textbf{(A)} RMS vs decay parameter $\lambda$ for different maximum lookback days. \textbf{(B)} Heatmap of RMS across all configurations. \textbf{(C)} Weight profiles showing how different $\lambda$ values distribute weight across lookback days. \textbf{(D)} Summary: best configuration (5 days, $\lambda$=0.5) achieves only 8 milli-degree improvement.}
\label{fig:exponential}
\end{figure}

\textbf{Key finding}: The best configuration (5 days lookback with $\lambda=0.5$) achieves 1.050°C RMS vs 1.058°C baseline---an improvement of only \textbf{8 milli-degrees}. Exponential weighting provides minimal benefit over uniform weighting with integral day boundaries.

\subsection{Hybrid Predictor Design}

We investigated whether condition-specific models could outperform a single unified model.

\begin{figure}[H]
\centering
\includegraphics[width=\textwidth]{fig28_hybrid_predictor.png}
\caption{Hybrid predictor analysis. \textbf{(A)} Comparison of single RF baseline, condition-specific models, and weighted ensemble. \textbf{(B)} Performance breakdown by stability condition. \textbf{(C)} Uncertainty calibration showing predicted vs actual errors. \textbf{(D)} Summary: condition-specific models provide marginal improvement; rate of change is best used for uncertainty quantification.}
\label{fig:hybrid}
\end{figure}

Three hybrid approaches were tested:
\begin{itemize}[noitemsep]
    \item \textbf{Condition-specific models}: Separate models for stable/unstable conditions $\rightarrow$ 1.089°C (worse)
    \item \textbf{Condition-weighted ensemble}: Weight models by stability $\rightarrow$ 1.064°C (+10 milli-degrees)
    \item Performance by condition: Stable = 0.99°C, Unstable = 1.11°C
\end{itemize}

\subsection{Rate-of-Change Weighted Inputs}

We tested using rate of change to adaptively weight lookback observations and training samples.

\begin{figure}[H]
\centering
\includegraphics[width=\textwidth]{fig29_rate_weighted_inputs.png}
\caption{Rate-weighted input analysis. \textbf{(A)} RMS vs rate weight factor---higher factors weight recent data more heavily when temperatures are changing rapidly. \textbf{(B)} Comparison of sample weighting schemes. \textbf{(C)} Trade-off between stable and unstable condition performance. \textbf{(D)} Summary: sample weighting with 3$\times$ focus on unstable conditions yields best improvement.}
\label{fig:rate_weighted}
\end{figure}

\begin{table}[H]
\centering
\caption{Rate-Weighted Input Results}
\label{tab:rate_weighted}
\begin{tabular}{lcccc}
\toprule
\textbf{Method} & \textbf{Overall RMS} & \textbf{Stable RMS} & \textbf{Unstable RMS} & \textbf{Improvement} \\
\midrule
Baseline (uniform) & 1.059°C & 0.999°C & 1.086°C & --- \\
Rate-weighted inputs ($\lambda$=1.0) & 1.052°C & 0.984°C & 1.082°C & +7 m°C \\
Sample weight (3$\times$ unstable) & 1.047°C & 0.987°C & 1.074°C & +12 m°C \\
\bottomrule
\end{tabular}
\end{table}

\textbf{Key finding}: Weighting training samples 3$\times$ for unstable conditions provides \textbf{+12 milli-degree improvement}. The primary value of rate-of-change analysis remains uncertainty quantification rather than dramatic accuracy gains.

\subsection{Summary of Additional Methodologies}

Table~\ref{tab:methodology_summary} summarizes all additional approaches tested.

\begin{table}[H]
\centering
\caption{Summary of Additional Methodology Exploration}
\label{tab:methodology_summary}
\begin{tabular}{lcc}
\toprule
\textbf{Methodology} & \textbf{Result} & \textbf{Recommendation} \\
\midrule
Exponential weighting & +8 m°C improvement & Marginal benefit \\
Condition-specific models & $-$15 m°C (worse) & Not recommended \\
Condition-weighted ensemble & +10 m°C improvement & Optional \\
Rate-weighted inputs & +7 m°C improvement & Marginal benefit \\
Sample weighting (3$\times$ unstable) & +12 m°C improvement & Recommended \\
\midrule
\textbf{Rate of change for uncertainty} & \textbf{Identifies 30\% error diff.} & \textbf{Highly recommended} \\
\bottomrule
\end{tabular}
\end{table}

\textbf{Main conclusion}: While various weighting schemes provide modest improvements (7--12 milli-degrees), the most valuable finding is that \textbf{rate of change reliably predicts when predictions will be less accurate}. This enables meaningful uncertainty quantification: predictions made during stable conditions (RMS $\approx$ 0.9°C) are more reliable than those made during rapidly changing conditions (RMS $\approx$ 1.1°C).

\subsection{Comprehensive Comparison: All Methods}

Figure~\ref{fig:comprehensive} provides a final visual summary comparing all numerical and text-based approaches across lead times.

\begin{figure}[H]
\centering
\includegraphics[width=\textwidth]{fig30_comprehensive_summary.png}
\caption{Comprehensive comparison of all prediction methods. \textbf{(A)} All methods ranked by RMS at 3-hour lead time (blue=numerical, coral=text-based, gray=baseline). \textbf{(B)} RMS vs lead time for key methods showing Random Forest dominates at all lead times. \textbf{(C)} Best numerical vs best text-based performance at each lead time. \textbf{(D)} Summary of key findings.}
\label{fig:comprehensive}
\end{figure}

Table~\ref{tab:all_leadtimes} presents performance across all lead times for the top methods.

\begin{table}[H]
\centering
\caption{Top Methods Across All Lead Times}
\label{tab:all_leadtimes}
\begin{tabular}{llcccc}
\toprule
\textbf{Lead Time} & \textbf{Method} & \textbf{RMS (°C)} & \textbf{MAE (°C)} & \textbf{Bias (°C)} & \textbf{\% $<$ 1°C} \\
\midrule
\multirow{4}{*}{1 hour}
& Ridge & 0.715 & 0.569 & $+$0.06 & 84.4 \\
& \textbf{Random Forest} & \textbf{0.719} & 0.568 & $+$0.07 & 85.4 \\
& Transformer & 0.87 & 0.69 & $+$0.05 & 75.9 \\
& LSTM+Attention & 0.93 & 0.72 & $+$0.12 & 74.5 \\
\midrule
\multirow{4}{*}{2 hours}
& \textbf{Random Forest} & \textbf{0.946} & 0.740 & $+$0.10 & 74.0 \\
& RF + Sample Weight & 0.946 & 0.739 & $+$0.11 & 72.9 \\
& LSTM+Attention & 1.14 & 0.91 & $+$0.19 & 59.4 \\
& Transformer & 1.22 & 0.99 & $+$0.19 & 58.4 \\
\midrule
\multirow{4}{*}{3 hours}
& \textbf{RF + Sample Weight} & \textbf{1.047} & 0.812 & $+$0.11 & 68.7 \\
& Random Forest & 1.057 & 0.822 & $+$0.12 & 69.0 \\
& LSTM+Attention & 1.31 & 1.04 & $+$0.41 & 54.9 \\
& Transformer & 1.43 & 1.15 & $+$0.25 & 52.3 \\
\bottomrule
\end{tabular}
\end{table}

\section{Key Findings}

\begin{enumerate}
    \item \textbf{61\% improvement over baseline}: Best ML model achieves 1.01°C RMS vs 2.57°C for persistence.

    \item \textbf{Near-zero bias}: Calibration-based correction reduces systematic bias to $<$0.05°C.

    \item \textbf{70\% of predictions within 1°C}: At 3 hours, most predictions are highly accurate.

    \item \textbf{Median error of 0.64°C}: Half of all predictions are within 0.64°C of actual.

    \item \textbf{72-hour lookback optimal}: Three days of historical data captures diurnal patterns without adding noise.

    \item \textbf{Tree-based models dominate}: Random Forest and gradient boosting variants outperform neural networks and linear models.

    \item \textbf{Error grows at 0.177°C/hour}: Fine-grained analysis shows smooth degradation from 0.30°C at 15 minutes to 1.07°C at 3 hours, crossing the 1.0°C threshold at approximately 2.75 hours before sunset.

    \item \textbf{Rate of change predicts uncertainty}: The 2-hour rate of temperature change identifies when predictions will be less reliable---stable conditions yield RMS=0.88°C vs unstable conditions at RMS=1.14°C (30\% difference).

    \item \textbf{Alternative methodologies provide marginal gains}: Exponential weighting (+8 m°C), hybrid predictors (+10 m°C), and sample weighting (+12 m°C) were tested; none dramatically improve upon the baseline Random Forest approach.
\end{enumerate}

\section{Recommendations}

\subsection{For Operational Deployment}
\begin{enumerate}[noitemsep]
    \item Use \textbf{Random Forest} with 3-day lookback features
    \item Implement \textbf{rolling bias correction} using recent prediction errors
    \item \textbf{Report uncertainty based on rate of change}: When $|$rate$|<$0.15°C/h, report $\pm$0.9°C confidence; otherwise report $\pm$1.1°C
    \item Consider \textbf{3$\times$ sample weighting} for unstable conditions during training for modest improvement
    \item Retrain monthly to incorporate new data
    \item Monitor for anomalous days (frontal passages, unusual weather)
\end{enumerate}

\subsection{Operational Trigger Points}
Based on fine-grained lead time analysis, predictions can be made at any time with these expected accuracies:
\begin{itemize}[noitemsep]
    \item \textbf{3 hours}: RMS $\approx$ 1.0°C, 67\% within 1°C
    \item \textbf{2 hours}: RMS $\approx$ 0.95°C, 74\% within 1°C
    \item \textbf{1 hour}: RMS $\approx$ 0.72°C, 87\% within 1°C
    \item \textbf{15 minutes}: RMS $\approx$ 0.30°C, 99\% within 1°C
\end{itemize}

\subsection{Potential Future Improvements}
\begin{enumerate}[noitemsep]
    \item Incorporate additional meteorological data (cloud cover, pressure, humidity)
    \item Implement formal uncertainty quantification using rate of change as a predictor
    \item Test ensemble methods with condition-dependent weighting
\end{enumerate}

\section{Conclusion}

We developed a sunset temperature prediction system achieving \textbf{1.01°C RMS error} at 3 hours before sunset with \textbf{bias of only $-$0.04°C}. Using odd calendar days for training and even days for testing ensures rigorous validation. The system achieves a 61\% improvement over the persistence baseline.

Extensive additional methodology exploration confirmed that Random Forest with 3-day lookback remains the optimal approach. Alternative methods (exponential weighting, hybrid predictors, sample weighting) provide at most 12 milli-degree improvement. The most valuable finding from this exploration is that the \textbf{2-hour rate of temperature change reliably predicts prediction uncertainty}, enabling meaningful confidence intervals for operational use. The system is ready for deployment at the Rubin Observatory.

\vspace{1em}
\hrule
\vspace{0.5em}
\small
\textit{Report generated January 2026. Code available in project repository.}

\end{document}
